% =====================================================================
%  TCET CODING PLATFORM
%  Official Technical Documentation -- Draft 1 (valid until Stage 2)
%  MODULE 1 of 7 : Introduction & System Overview
%
%  STANDALONE FILE. Compile on its own with:
%     pdflatex Module1_Introduction.tex   (run twice for the TOC)
%  Powered by the Centre of Excellence (CoE), TCET.
% =====================================================================
\documentclass[12pt,a4paper,oneside]{report}

% ---------------------------------------------------------------------
%  Packages (self-contained; identical preamble across all 7 modules)
% ---------------------------------------------------------------------
\usepackage[T1]{fontenc}
\usepackage[utf8]{inputenc}
\usepackage[a4paper,margin=1in]{geometry}
\usepackage{mathptmx}
\usepackage{microtype}
\usepackage{graphicx}
\usepackage{xcolor}
\usepackage{tabularx}
\usepackage{booktabs}
\usepackage{longtable}
\usepackage{xltabular}
\usepackage{array}
\usepackage{enumitem}
\usepackage{listings}
\usepackage{titlesec}
\usepackage{fancyhdr}
\usepackage{parskip}
\usepackage[hidelinks]{hyperref}

\definecolor{tcetblue}{RGB}{0,0,0}
\definecolor{tcetaccent}{RGB}{0,0,0}
\definecolor{codebg}{RGB}{245,246,248}
\definecolor{codeframe}{RGB}{215,219,225}
\definecolor{codekw}{RGB}{0,0,0}
\definecolor{codecomment}{RGB}{0,0,0}
\definecolor{codestr}{RGB}{0,0,0}
\definecolor{boxbg}{RGB}{239,246,250}
\definecolor{boxframe}{RGB}{0,0,0}

\hypersetup{
  colorlinks=true, linkcolor=tcetblue, urlcolor=tcetaccent, citecolor=tcetblue,
  pdftitle={TCET Coding Platform -- Module 1: Introduction \& System Overview},
  pdfauthor={Centre of Excellence, TCET}
}

\lstdefinestyle{tcp}{
  backgroundcolor=\color{codebg}, basicstyle=\ttfamily\footnotesize,
  keywordstyle=\color{codekw}\bfseries, commentstyle=\color{codecomment}\itshape,
  stringstyle=\color{codestr}, numberstyle=\tiny\color{gray},
  frame=single, rulecolor=\color{codeframe}, breaklines=true,
  breakatwhitespace=false, showstringspaces=false, tabsize=2,
  columns=fullflexible, keepspaces=true, xleftmargin=0.6em, framexleftmargin=0.6em,
}
\lstset{style=tcp}

\newcommand{\code}[1]{\texttt{#1}}

% Lightweight callout box (no external packages)
\newenvironment{callout}[1]{%
  \par\vspace{4pt}\noindent
  \begin{tabular}{@{}p{0.014\linewidth}@{\hspace{0.6em}}p{0.93\linewidth}@{}}
  \textcolor{boxframe}{\rule{2pt}{\baselineskip}} &
  \textbf{\textcolor{tcetblue}{#1}}\\[2pt]
  & }{\end{tabular}\par\vspace{6pt}}

\titleformat{\chapter}[display]{\bfseries\fontsize{18}{22}\selectfont}{\bfseries\fontsize{14}{17}\selectfont Chapter \thechapter}{0.2em}{}
\titlespacing*{\chapter}{0pt}{6pt}{22pt}
\titleformat{\section}{\bfseries\fontsize{16}{19}\selectfont}{\thesection}{0.6em}{}
\titleformat{\subsection}{\bfseries\fontsize{14}{17}\selectfont}{\thesubsection}{0.5em}{}

\pagestyle{fancy}
\fancyhf{}
\renewcommand{\headrulewidth}{0.4pt}
\fancyhead[L]{\small\itshape TCET Coding Platform}
\fancyhead[R]{\small\itshape Module 1 \textbullet\ Introduction}
\fancyfoot[C]{\thepage}
\fancypagestyle{plain}{\fancyhf{}\fancyfoot[C]{\thepage}\renewcommand{\headrulewidth}{0pt}}

\newcommand{\note}[1]{\par\vspace{3pt}\noindent{\color{tcetaccent}\textbf{Note.}}~#1\par\vspace{3pt}}

% =====================================================================
\begin{document}

% ----------------------------- TITLE ---------------------------------
\begin{titlepage}
\centering
\vspace*{1.4cm}
{\color{tcetaccent}\rule{\linewidth}{1.2pt}}\\[0.5cm]
{\LARGE\bfseries TCET Coding Platform\par}
\vspace{0.4cm}
{\Large Official Technical Documentation\par}
{\large Draft~1 \textemdash\ Testing Phase\par}
{\color{tcetaccent}\rule{\linewidth}{1.2pt}}\\[1.4cm]
{\Huge\bfseries\color{tcetblue} Module 1\par}
\vspace{0.3cm}
{\Huge Introduction \&\\[0.2cm] System Overview\par}
\vspace{2.2cm}
{\large\bfseries Powered by\par}
{\Large\color{tcetblue} Centre of Excellence (CoE), TCET\par}
\vfill
\begin{tabular}{rl}
\textbf{Module} & 1 of 7\\[3pt]
\textbf{Scope} & Vision, scope, stakeholders, requirements, glossary\\[3pt]
\textbf{Status} & Draft~1 (living document until Stage~2)\\[3pt]
\textbf{Audience} & TCET Officials, CoE Review Committee, Engineering\\
\end{tabular}
\vspace{0.8cm}
{\small\itshape \today\par}
\end{titlepage}

\pagenumbering{roman}
\tableofcontents
\clearpage
\pagenumbering{arabic}

% =====================================================================
\chapter{Document Purpose \& Context}
% =====================================================================

\section{About this Module}
This is the first of seven modules that together constitute the official
technical documentation of the \textbf{TCET Coding Platform}. The platform is
developed and operated under
the \textbf{Centre of Excellence (CoE)} at Thakur College of Engineering and
Technology (TCET). This module establishes the shared context that the remaining
six modules build upon: what the platform is, why it exists, who it serves, the
boundaries of the current build, the requirements it is held to, and the
vocabulary used throughout the documentation set.

The documentation is deliberately partitioned into seven self-contained modules
so that distinct concerns can be reviewed, approved, versioned, and extended
independently by different readers \textemdash\ officials, security reviewers,
and engineers \textemdash\ without any one of them needing to read the entire
corpus to obtain the information relevant to their role.

\begin{callout}{Document set at a glance}
Module~1 Introduction \& System Overview \textbullet\ Module~2 System
Architecture \& Technology Stack \textbullet\ Module~3 Authentication,
Authorization, Security \& Data Persistence \textbullet\ Module~4 Problems,
Submissions \& Code Execution Engine \textbullet\ Module~5 Contest \& Assessment
Engine \textbullet\ Module~6 Proctoring \& Academic Integrity \textbullet\
Module~7 Deployment, Infrastructure, Operations \& Future Scope.
\end{callout}

\section{Status of this Document}
This is \textbf{Draft~1}. It is an authoritative but living record: it captures
the platform as it exists at the time of writing and is expected to be revised
at the Stage~2 review. Where this document and older, ad-hoc notes or README
files disagree, this document is the governing reference. Any figure, limit, or
behaviour stated here reflects the implemented system and is traceable to the
codebase; forward-looking items are explicitly confined to the Future Scope
material in Module~7 and are labelled as candidate, not committed, work.

\section{Intended Audience}
The documentation is written for three overlapping audiences:

\begin{itemize}[leftmargin=1.5em]
  \item \textbf{TCET Officials and the CoE Review Committee}, who require a clear,
        non-implementation view of what the platform does, how it protects
        examination integrity, how it scales to institutional load, and how it
        is operated.
  \item \textbf{Security and compliance reviewers}, who require precise detail on
        identity federation, the trust boundary, authorization, data handling,
        and the academic-integrity controls.
  \item \textbf{Engineers}, who require the internal contracts, data models,
        execution pipeline, and operational procedures needed to maintain and
        extend the platform.
\end{itemize}

\section{How to Read this Document Set}
\begin{longtable}{@{}p{0.32\linewidth} p{0.60\linewidth}@{}}
\toprule
\textbf{If you are\ldots} & \textbf{Read\ldots}\\
\midrule
\endhead
An official seeking an overview & Module~1 in full; the introductions of Modules~2, 5, 6, and 7.\\
Assessing examination integrity & Module~5 (contest engine) and Module~6 (proctoring), plus the security controls in Module~3.\\
Assessing security posture & Module~3 in full, plus Module~2's trust-boundary discussion.\\
Assessing scalability / operations & Module~2 (architecture) and Module~7 (deployment, scaling, runbooks).\\
Engineering / maintaining the code & All modules; Modules~3--6 carry the internal contracts and data models.\\
\bottomrule
\end{longtable}

\section{Document Conventions}
\begin{itemize}[leftmargin=1.5em]
  \item Fixed-width text such as \code{x-coe-email} or \code{contest\_attempts}
        denotes an identifier, header, field name, endpoint, environment
        variable, or file path taken from the implementation.
  \item Roles are written in small capitals of meaning: a \emph{Student} and a
        \emph{Faculty} member are the two internal roles; the upstream identity
        provider additionally issues \emph{Admin} and \emph{Industry} roles that
        the platform maps onto Faculty.
  \item ``The platform'' and ``the system'' are used
        interchangeably.
  \item Time-relative statements are avoided; where a date matters it is stated
        absolutely.
\end{itemize}

% =====================================================================
\chapter{Product Vision \& Objectives}
% =====================================================================

\section{The Problem being Solved}
Engineering education at scale requires frequent, hands-on programming practice
and assessment. Conducting these activities with general-purpose tools
(spreadsheets for marks, email or shared drives for submissions, individual
laboratory machines for compilation) does not scale to an institution running
mass laboratory sessions, and it provides weak guarantees of academic integrity
during timed examinations. Three structural problems recur:

\begin{enumerate}[leftmargin=1.6em]
  \item \textbf{Throughput.} When thousands of students submit code within the
        same short window, any system that compiles and runs code
        \emph{synchronously} on the request path collapses under the load.
  \item \textbf{Integrity.} Timed coding examinations conducted in an ordinary
        browser are trivially undermined by tab-switching, copy/paste, external
        lookups, and screenshots, and by any design that reveals correctness
        during the examination.
  \item \textbf{Fragmentation and ownership.} Identity, problem banks, results,
        and analytics scattered across third-party tools are hard to govern,
        hard to audit, and not owned by the institution.
\end{enumerate}

\section{The Vision}
The TCET Coding Platform is a single, institution-owned environment in which faculty author
programming problems and timed contests, and students solve them with immediate,
sandboxed code execution \textemdash\ engineered from the outset for mass
laboratory use, strong examination integrity, and institutional ownership of
identity and data. The platform is built on three guiding principles.

\begin{callout}{Guiding principles}
\textbf{Scale} \textemdash\ absorb bursts of thousands of concurrent submissions
by decoupling request handling from code execution through an asynchronous
queue. \quad \textbf{Integrity} \textemdash\ enforce a browser-level proctoring
regime and a deferred-grading model so no student can learn or leak answers
during a live examination. \quad \textbf{Institutional ownership} \textemdash\
federate identity through the CoE's centralized Single Sign-On and never expose
the backend directly to the public internet.
\end{callout}

\section{Objectives \& Success Criteria}
\begin{longtable}{@{}p{0.30\linewidth} p{0.62\linewidth}@{}}
\toprule
\textbf{Objective} & \textbf{Success criterion}\\
\midrule
\endhead
Mass-scale practice \& assessment & The system sustains large concurrent
submission bursts without loss of submissions or request-path failure, by
buffering work on a queue and draining it at worker capacity.\\
Trustworthy examinations & During a live contest no score, correctness, or
correct answer is revealed to a student; integrity violations are detected,
recorded, and penalized; abandoned attempts are still graded fairly.\\
Institution-owned identity & All authentication is federated from the CoE SSO;
the platform stores no passwords and exposes no public authentication surface.\\
Multi-language reach & Students may solve in 20+ programming languages with a
consistent editor and judging experience.\\
Actionable analytics & Faculty and students obtain ratings, accuracy, activity
heatmaps, standings, and exportable results.\\
Operational resilience & The platform self-heals from worker crashes, queue
backpressure, and students disconnecting mid-contest.\\
\bottomrule
\end{longtable}

\section{Design Values}
Beyond the headline objectives, the platform embodies a set of engineering
values that recur throughout the later modules and explain many of its design
choices:

\begin{description}[leftmargin=1.6em,style=nextline]
  \item[Defence in depth] Security and correctness are enforced at multiple
        independent layers (proxy trust, identity verification, role checks,
        ownership checks, schema validation), so no single failure is
        catastrophic.
  \item[Determinism] Time-dependent and scoring logic is written to be
        deterministic and idempotent, so the same inputs always produce the same
        result and operations can be safely retried.
  \item[Fail-safe over fail-open] When something goes wrong (a queue failure, a
        judging error, a disconnect), the system defaults to the safe outcome
        (mark the submission errored, finalize the attempt) rather than leaving
        state ambiguous.
  \item[Snapshots over joins] Historical records embed the values they depend on
        at the time they were created, so history stays stable and reads stay
        cheap.
  \item[Testability] Every layer depends on interfaces and an injected clock, so
        the whole system can be exercised deterministically without production
        infrastructure.
\end{description}

% =====================================================================
\chapter{Scope}
% =====================================================================

\section{In Scope (current build)}
The following capabilities are implemented and are documented across Modules~2--7:

\begin{itemize}[leftmargin=1.5em]
  \item Role-based access for \textbf{Students} and \textbf{Faculty}, federated
        via the CoE SSO, including profile completion.
  \item A \textbf{problem bank} with lifecycle management (Draft / Published /
        Archived), difficulty tiers, tags, per-department targeting, resource
        limits, and both sample and hidden test cases.
  \item An asynchronous \textbf{submission and judging pipeline} backed by a
        Redis queue and Judge0 sandboxed execution, supporting 20+ executable
        languages and three editor-only languages.
  \item A full \textbf{contest and assessment engine} supporting MCQ, MSQ, and
        Coding questions, per-student attempt windows, deferred grading,
        automatic finalization, standings, faculty review, and CSV export.
  \item A \textbf{proctoring and academic-integrity subsystem} for live contests.
  \item A \textbf{post-contest feedback} mechanism that gates the release of a
        student's published results behind a usability questionnaire.
  \item \textbf{Analytics}: per-user ratings, accuracy, submission heatmaps,
        difficulty and language breakdowns, and leaderboards.
\end{itemize}

\section{Out of Scope (this draft)}
\begin{itemize}[leftmargin=1.5em]
  \item The CoE SSO identity provider itself, which is owned and operated
        upstream; the platform only consumes its signed identity.
  \item Public or unauthenticated access; the backend is designed to sit behind a
        trusted reverse proxy at all times.
  \item Native mobile applications; the frontend is a responsive web application.
  \item The forward-looking learning environments (AI-assisted analysis,
        practical labs, DBMS/SQL, Linux terminal) described as Future Scope in
        Module~7, which are candidate work, not part of the current build.
\end{itemize}

\section{Assumptions}
\begin{longtable}{@{}p{0.05\linewidth} p{0.87\linewidth}@{}}
\toprule
\textbf{\#} & \textbf{Assumption}\\
\midrule
\endhead
A1 & A trusted reverse proxy (Cloudflare Tunnel) terminates CoE authentication upstream and forwards a signed identity; the backend is never exposed directly to the public internet.\\
A2 & The reverse proxy strips any client-supplied \code{x-coe-*} headers so identity cannot be spoofed by clients.\\
A3 & MongoDB, Redis, and a Judge0 execution environment are reachable from the backend host.\\
A4 & Students take examinations in a modern desktop browser that supports the Fullscreen and Clipboard APIs used by the proctoring subsystem.\\
A5 & The CoE SSO issues one of the recognized roles (Admin, Faculty, Industry, Student) and an account status.\\
\bottomrule
\end{longtable}

\section{Constraints}
\begin{longtable}{@{}p{0.05\linewidth} p{0.87\linewidth}@{}}
\toprule
\textbf{\#} & \textbf{Constraint}\\
\midrule
\endhead
C1 & The platform does not manage passwords or credentials; it must rely entirely on the CoE SSO for authentication.\\
C2 & Untrusted student code must be executed only inside a sandbox; it must never run with network access or outside resource limits.\\
C3 & Request bodies are capped (100\,kb) and API calls are rate limited, so clients must operate within those limits.\\
C4 & Browser-based proctoring cannot control the operating system or external devices; integrity guarantees are bounded accordingly (see Module~6).\\
C5 & Editor-only languages (HTML, CSS, React) are never executed by the judge and exist for authoring/preview contexts only.\\
\bottomrule
\end{longtable}

\section{External Dependencies}
\begin{longtable}{@{}p{0.24\linewidth} p{0.68\linewidth}@{}}
\toprule
\textbf{Dependency} & \textbf{Role and consequence if unavailable}\\
\midrule
\endhead
CoE SSO & Federated identity. Without it, users cannot authenticate; the platform has no fallback credential store by design.\\
Reverse proxy (Cloudflare Tunnel) & Public routing and identity injection. Without it, the backend has no safe public entry point.\\
MongoDB & System of record. Without it, no persistence; the platform cannot serve.\\
Redis & Submission queue and buffering. Without it, submissions cannot be enqueued and are failed fast.\\
Judge0 & Sandboxed code execution. Without it, code cannot be run or judged; runs return an execution error.\\
\bottomrule
\end{longtable}

% =====================================================================
\chapter{Stakeholders \& Actors}
% =====================================================================

\section{Stakeholder Register}
\begin{longtable}{@{}p{0.24\linewidth} p{0.32\linewidth} p{0.32\linewidth}@{}}
\toprule
\textbf{Stakeholder} & \textbf{Primary interest} & \textbf{Interaction with the platform}\\
\midrule
\endhead
Centre of Excellence (CoE) & Ownership, identity federation, governance. & Operates the platform and the upstream SSO.\\
TCET Officials & Governance, approval, Stage-gate review. & Review documentation; authorize progression.\\
Faculty & Authoring, assessment, integrity, review. & Create problems and contests; proctor; review and export results.\\
Students & Practice and assessment. & Solve problems; register for and attempt contests; view analytics.\\
Platform / DevOps team & Reliability, scaling, operations. & Deploy, monitor, scale, and maintain the system.\\
Security reviewers & Confidentiality, integrity, auditability. & Assess the trust boundary and integrity controls.\\
\bottomrule
\end{longtable}

\section{System Actors}
An \emph{actor} is any entity that interacts with the platform's boundary.

\begin{longtable}{@{}p{0.24\linewidth} p{0.68\linewidth}@{}}
\toprule
\textbf{Actor} & \textbf{Description}\\
\midrule
\endhead
Student (human) & Authenticated learner who solves problems and attempts contests under the Student role.\\
Faculty (human) & Authenticated educator who authors and administers content under the Faculty role (also the mapping target for Admin and Industry).\\
CoE SSO (system) & Upstream identity provider that authenticates users and issues signed identity.\\
Reverse proxy (system) & Trusted intermediary that forwards authenticated requests with identity headers.\\
Judge0 (system) & Sandboxed execution engine invoked by the submission worker.\\
Submission worker (system) & Background process that consumes the queue and drives judging.\\
Background schedulers (system) & In-process timers that recover stale submissions and finalize expired attempts.\\
\bottomrule
\end{longtable}

\section{Role Capability Matrix}
The table summarizes, at a glance, what each human role may do. Precise
authorization rules are in Module~3.

\begin{longtable}{@{}p{0.50\linewidth} c c@{}}
\toprule
\textbf{Capability} & \textbf{Student} & \textbf{Faculty}\\
\midrule
\endhead
Solve published problems (run / submit) & Yes & --\\
View own submissions and analytics & Yes & Yes\\
Register for and attempt contests & Yes & --\\
Submit post-contest feedback & Yes & --\\
Author / edit / archive problems & -- & Yes\\
Create / edit contests and questions & -- & Yes\\
View any submission for owned resources & -- & Yes\\
Review contest attempts (incl. final code) & -- & Yes\\
Publish contest results & -- & Yes\\
Export standings / registrations / leaderboard (CSV) & -- & Yes\\
View another student's profile analytics & -- & Yes\\
\bottomrule
\end{longtable}

% =====================================================================
\chapter{Functional Capability Catalogue}
% =====================================================================
This chapter enumerates the platform's functional capabilities at a level
suitable for officials, with a one-line pointer to the module that specifies
each in depth.

\section{Identity \& Access}
\begin{longtable}{@{}p{0.26\linewidth} p{0.52\linewidth} p{0.12\linewidth}@{}}
\toprule
\textbf{Capability} & \textbf{Description} & \textbf{Module}\\
\midrule
\endhead
Federated login & Authenticate via CoE SSO using signed identity headers or a JWT; establish a Student or Faculty session. & 3\\
Profile completion & Capture department, semester, roll number, and profile links required for eligibility and analytics. & 3\\
Role-based routing & Serve distinct Student and Faculty application areas gated by role. & 2, 3\\
Session termination & Clear federated cookies and return to an allowlisted frontend origin. & 3\\
\bottomrule
\end{longtable}

\section{Problem Bank \& Practice}
\begin{longtable}{@{}p{0.26\linewidth} p{0.52\linewidth} p{0.12\linewidth}@{}}
\toprule
\textbf{Capability} & \textbf{Description} & \textbf{Module}\\
\midrule
\endhead
Problem authoring & Rich statements, constraints, formats, tags, difficulty, resource limits, sample \& hidden tests. & 4\\
Lifecycle control & Move problems through Draft, Published, and Archived states. & 4\\
Department targeting & Restrict a problem's visibility to a department, or leave it open to all. & 4\\
Run (samples) & Execute code against sample tests for immediate, unscored feedback. & 4\\
Submit (all tests) & Queue a submission judged against all tests; award first-solve rating. & 4\\
Auto-wrapping & Accept LeetCode-style Solution classes and synthesize a runnable entry point. & 4\\
\bottomrule
\end{longtable}

\section{Contests \& Assessment}
\begin{longtable}{@{}p{0.26\linewidth} p{0.52\linewidth} p{0.12\linewidth}@{}}
\toprule
\textbf{Capability} & \textbf{Description} & \textbf{Module}\\
\midrule
\endhead
Contest authoring & Define window, per-attempt duration, registration window, type, targeting, violation limit, and questions. & 5\\
Question types & MCQ, MSQ, and Coding questions in a single contest. & 5\\
Registration \& attempts & Register within the window; start a single attempt with a frozen personal deadline. & 5\\
Deferred grading & Score nothing live; grade at submit, auto-submit, and publish. & 5\\
Auto-submission & Submit unsubmitted drafts when an attempt ends by any route. & 5\\
Standings \& review & Ranked standings; faculty attempt review including final code. & 5\\
Feedback gate & Require a usability questionnaire before releasing a student's published results. & 5\\
\bottomrule
\end{longtable}

\section{Integrity \& Analytics}
\begin{longtable}{@{}p{0.26\linewidth} p{0.52\linewidth} p{0.12\linewidth}@{}}
\toprule
\textbf{Capability} & \textbf{Description} & \textbf{Module}\\
\midrule
\endhead
Proctoring enforcement & Fullscreen lock, focus / tab / screenshot detection, shortcut and clipboard blocking, watermark, screen guard. & 6\\
Violation scoring & Penalize unpreventable evasions; auto-submit at the violation threshold. & 5, 6\\
Event logging & Persist an immutable log of proctoring events for faculty review. & 6\\
Ratings \& accuracy & Difficulty-weighted first-solve rating, accuracy, problems solved. & 4\\
Leaderboards & Denormalized ranking of users. & 4\\
Profile analytics & Difficulty and language breakdowns, activity heatmap, recent activity. & 4\\
\bottomrule
\end{longtable}

% =====================================================================
\chapter{Non-Functional Requirements}
% =====================================================================
Non-functional requirements (NFRs) express the qualities the platform must
exhibit. Each is stated with the mechanism that satisfies it and the module that
elaborates it.

\section{Performance \& Scalability}
\begin{longtable}{@{}p{0.10\linewidth} p{0.50\linewidth} p{0.28\linewidth}@{}}
\toprule
\textbf{Ref} & \textbf{Requirement} & \textbf{Mechanism}\\
\midrule
\endhead
NFR-P1 & Sustain large concurrent submission bursts without dropping work. & Redis/BullMQ queue buffers submissions; workers drain at capacity (M2, M4).\\
NFR-P2 & Keep request latency independent of judging latency. & Submit is asynchronous; a dedicated worker process handles execution (M2, M7).\\
NFR-P3 & Judge fairly across fast and slow runtimes. & Per-language CPU-time and memory multipliers, capped by a ceiling (M4).\\
NFR-P4 & Scale judging throughput horizontally. & Multiple worker processes/hosts against one queue; tunable concurrency (M7).\\
\bottomrule
\end{longtable}

\section{Security}
\begin{longtable}{@{}p{0.10\linewidth} p{0.50\linewidth} p{0.28\linewidth}@{}}
\toprule
\textbf{Ref} & \textbf{Requirement} & \textbf{Mechanism}\\
\midrule
\endhead
NFR-S1 & Only authenticated, trusted-proxy traffic may reach the API. & Trusted-proxy source-IP allowlist; identity verification (M3).\\
NFR-S2 & No credential storage. & Federated identity via CoE SSO (M3).\\
NFR-S3 & Prevent cross-site request forgery on mutations. & Mutation-origin guard against an allowlist (M3).\\
NFR-S4 & Contain untrusted code execution. & Judge0 \code{isolate} sandbox, no network, resource-limited (M4).\\
NFR-S5 & Protect against request floods. & Global and per-user rate limiting; payload cap (M3).\\
\bottomrule
\end{longtable}

\section{Reliability \& Availability}
\begin{longtable}{@{}p{0.10\linewidth} p{0.50\linewidth} p{0.28\linewidth}@{}}
\toprule
\textbf{Ref} & \textbf{Requirement} & \textbf{Mechanism}\\
\midrule
\endhead
NFR-R1 & No submission is permanently stranded by a worker crash. & Stale-submission recovery re-enqueues stuck work (M2, M4).\\
NFR-R2 & Abandoned attempts are graded near their real deadline. & Background attempt finalizer (M2, M5).\\
NFR-R3 & Operations are safe to retry. & Idempotent finalization and scoring (M4, M5).\\
NFR-R4 & Processes recover from failure automatically. & PM2 autorestart with memory ceilings (M7).\\
\bottomrule
\end{longtable}

\section{Integrity, Usability, Maintainability \& Portability}
\begin{longtable}{@{}p{0.10\linewidth} p{0.50\linewidth} p{0.28\linewidth}@{}}
\toprule
\textbf{Ref} & \textbf{Requirement} & \textbf{Mechanism}\\
\midrule
\endhead
NFR-I1 & No answer leakage during a live contest. & Deferred grading; results revealed only at publish (M5).\\
NFR-I2 & Detect and penalize examination evasion. & Browser proctoring with server-side scoring (M6).\\
NFR-U1 & Consistent, accessible, themeable UI. & Tailwind + accessible component library; light/dark themes (M2).\\
NFR-U2 & Immediate feedback while coding. & In-browser editor with Run mode (M2, M4).\\
NFR-M1 & Storage-agnostic, testable business logic. & Layered architecture; repository interfaces; injected clock (M2).\\
NFR-M2 & Fail-fast on misconfiguration. & Startup-time schema validation of configuration (M2, M7).\\
NFR-Po1 & Portable execution across deployments. & Pluggable execution provider (Judge0 or stub) (M4).\\
\bottomrule
\end{longtable}

% =====================================================================
\chapter{Representative Use Cases}
% =====================================================================
The following use cases illustrate the principal end-to-end flows. Each names the
actor, preconditions, the main flow, and notable alternate outcomes. Detailed
behaviour is specified in the referenced modules.

\section{UC-1: Student Solves a Practice Problem}
\textbf{Actor:} Student. \quad \textbf{Precondition:} Authenticated; problem is
Published and visible to the student's department.
\begin{enumerate}[leftmargin=1.6em]
  \item The student opens a problem and writes a solution in the editor.
  \item The student presses \emph{Run}; the code executes against sample tests
        and returns immediate, unscored feedback.
  \item The student presses \emph{Submit}; a submission is persisted and queued,
        and a receipt is returned.
  \item A worker judges the submission against all tests; the verdict is stored.
  \item On first acceptance, difficulty-weighted rating is awarded and the
        student's analytics and the leaderboard are updated.
\end{enumerate}
\textbf{Alternates:} enqueue failure marks the submission errored immediately; a
non-accepted verdict awards no rating; re-solving awards nothing further.

\section{UC-2: Faculty Authors and Publishes a Problem}
\textbf{Actor:} Faculty. \quad \textbf{Precondition:} Authenticated as Faculty.
\begin{enumerate}[leftmargin=1.6em]
  \item The faculty member creates a problem in Draft with statement,
        constraints, limits, and sample and hidden tests.
  \item The problem is optionally targeted to a department.
  \item The faculty member publishes the problem, making it visible to eligible
        students; it may later be archived.
\end{enumerate}

\section{UC-3: Student Takes a Proctored Contest}
\textbf{Actor:} Student. \quad \textbf{Precondition:} Registered within the
registration window; contest is Live.
\begin{enumerate}[leftmargin=1.6em]
  \item The student starts an attempt; a personal deadline is frozen as the
        lesser of the per-attempt duration and the contest end.
  \item The proctoring subsystem enforces fullscreen, blanks the contest on focus
        loss, and records violations.
  \item The student answers objective questions and writes/submits code; nothing
        is scored, and drafts are auto-saved.
  \item The attempt ends by manual submit, deadline expiry, or reaching the
        violation limit; unsubmitted drafts are auto-submitted.
  \item After faculty publishes results, the student submits the feedback
        questionnaire and then views their report and standings.
\end{enumerate}
\textbf{Alternates:} disconnecting is handled by the background finalizer, which
grades the attempt at its deadline.

\section{UC-4: Faculty Runs and Reviews a Contest}
\textbf{Actor:} Faculty.
\begin{enumerate}[leftmargin=1.6em]
  \item The faculty member creates a contest with a window, duration,
        registration window, questions, and violation limit.
  \item Students register and attempt within the window.
  \item After the contest ends, the faculty member publishes results, triggering
        an authoritative re-grade of every attempt.
  \item The faculty member reviews individual attempts (including final code) and
        exports standings as CSV.
\end{enumerate}

% =====================================================================
\chapter{Technology Overview \& Deployment in Brief}
% =====================================================================

\section{Technology Stack}
A detailed treatment with rationale appears in Module~2; the overview orients the
reader.

\begin{longtable}{@{}p{0.22\linewidth} p{0.70\linewidth}@{}}
\toprule
\textbf{Layer} & \textbf{Technology}\\
\midrule
\endhead
Frontend & React~18, Vite~5, TypeScript, React Router, TanStack Query, Tailwind CSS, an accessible component library, and the Monaco editor.\\
Backend & Node.js, Express~5, TypeScript, Zod validation, Helmet, rate limiting.\\
Data store & MongoDB (official driver) with server-side JSON Schema validation.\\
Queue \& cache & Redis with BullMQ for the submission queue.\\
Execution & Judge0 (sandboxed via \code{isolate}) with an optional cloud fallback and a stub provider for tests.\\
Identity & CoE centralized SSO consumed through trusted proxy headers / JWT.\\
Infrastructure & Cloudflare Tunnels, Tailscale, PM2, and Docker (for Judge0).\\
\bottomrule
\end{longtable}

\begin{callout}{On persistence naming}
Some backend repository classes carry a \code{Firestore}-prefixed name for
historical reasons, and an early note referenced ``Firebase / Firestore.'' The
authoritative and only datastore is \textbf{MongoDB}; the prefix is a legacy
artefact and does not indicate a second datastore. This is stated explicitly so
the naming is never mistaken for a Firestore dependency.
\end{callout}

\section{Deployment Model in Brief}
Production places the backend behind a trusted reverse proxy (Cloudflare Tunnel),
with the CoE SSO terminating authentication upstream and forwarding a signed
identity. Administrative access to the host is restricted to a Tailscale network.
Application processes (API, submission worker, frontend) are supervised by PM2,
and the Judge0 execution stack runs as a Docker deployment. Full topology,
configuration, and runbooks are in Module~7.

\section{Logical Component Summary}
\begin{lstlisting}[language={}]
Browser (SPA)
   |  HTTPS
   v
Reverse proxy / CoE SSO edge   (adds signed x-coe-* identity)
   |  trusted source IP
   v
Express API  --(persist)-->  MongoDB
   |  (enqueue submissions)
   v
Redis queue  -->  Submission worker  -->  Judge0 (isolate sandbox)
                        |
                        v
                 verdict persisted; aggregates + leaderboard updated
\end{lstlisting}

% =====================================================================
\chapter{Risks, Assumptions \& Mitigations}
% =====================================================================
This risk register records the principal risks to the platform's objectives, a
qualitative severity, and the mitigation in place or planned.

\begin{longtable}{@{}p{0.05\linewidth} p{0.42\linewidth} c p{0.30\linewidth}@{}}
\toprule
\textbf{ID} & \textbf{Risk} & \textbf{Sev.} & \textbf{Mitigation}\\
\midrule
\endhead
RK1 & Backend exposed directly to the internet, allowing header spoofing. & High & Trusted-proxy source-IP allowlist; proxy strips client \code{x-coe-*} (M3).\\
RK2 & Submission spike overwhelms execution. & High & Asynchronous queue buffering; horizontal workers (M2, M4, M7).\\
RK3 & Untrusted code harms the host. & High & \code{isolate} sandbox; no network; resource limits (M4).\\
RK4 & Answer leakage during examinations. & High & Deferred grading; results withheld until publish (M5).\\
RK5 & Examination evasion via tab-switch / screenshot. & Medium & Browser proctoring with server-side scoring; auto-submit at threshold (M6).\\
RK6 & Student disconnects mid-contest, losing work. & Medium & Draft auto-save and auto-submission; background finalizer (M5).\\
RK7 & Worker crash strands submissions. & Medium & Stale-submission recovery (M4).\\
RK8 & Misconfiguration starts the system in an insecure state. & Medium & Startup-time configuration validation; fail fast (M2, M7).\\
RK9 & Browser proctoring cannot stop external cameras / second devices. & Accepted & Watermark and screen guard deter; invigilation and future integrations (M6, M7).\\
\bottomrule
\end{longtable}

% =====================================================================
\chapter{Glossary \& Abbreviations}
% =====================================================================

\section{Terms}
\begin{longtable}{@{}p{0.26\linewidth} p{0.66\linewidth}@{}}
\toprule
\textbf{Term} & \textbf{Definition}\\
\midrule
\endhead
Attempt & A single student's run of a contest, with its own frozen deadline and per-question state.\\
Contest & A timed assessment of MCQ / MSQ / Coding questions authored by faculty.\\
Deferred grading & The model in which nothing is scored during a live contest; grading happens at submit, auto-submit, and publish.\\
Editor-only language & A language (HTML, CSS, React) that is never executed by the judge.\\
Feedback gate & The requirement that a student submit a usability questionnaire before their published contest results are released.\\
Finalization & The one-time application of a judged verdict to a submission, and the grading of an attempt.\\
Hidden test case & A test case used for judging but never shown to students.\\
Judge0 & The sandboxed, multi-language code-execution engine.\\
Lifecycle state & The Draft / Published / Archived status of a problem.\\
Proctoring event & A recorded integrity event raised during an attempt.\\
Rating & A difficulty-weighted score awarded for the first acceptance of a problem.\\
Sample test case & A visible test case shown to students and used by Run.\\
Snapshot & A denormalized copy of a value stored on a record to keep history stable.\\
Trusted proxy & The reverse proxy whose source IP the backend trusts to have authenticated the request.\\
\bottomrule
\end{longtable}

\section{Abbreviations}
\begin{longtable}{@{}p{0.18\linewidth} p{0.74\linewidth}@{}}
\toprule
\textbf{Abbrev.} & \textbf{Expansion}\\
\midrule
\endhead
API & Application Programming Interface.\\
CoE & Centre of Excellence (TCET).\\
CORS & Cross-Origin Resource Sharing.\\
CSRF & Cross-Site Request Forgery.\\
CSV & Comma-Separated Values.\\
JWT & JSON Web Token.\\
MCQ / MSQ & Multiple-Choice / Multiple-Select Question.\\
NFR & Non-Functional Requirement.\\
SPA & Single-Page Application.\\
SSO & Single Sign-On.\\
TCET & Thakur College of Engineering and Technology.\\
TLE & Time-Limit Exceeded.\\
UC & Use Case.\\
\bottomrule
\end{longtable}

% =====================================================================
\chapter{Module Roadmap}
% =====================================================================
The remaining modules develop the summary given here into full specifications.

\begin{longtable}{@{}p{0.06\linewidth} p{0.30\linewidth} p{0.54\linewidth}@{}}
\toprule
\textbf{Mod.} & \textbf{Title} & \textbf{Focus}\\
\midrule
\endhead
2 & System Architecture \& Technology Stack & Layered architecture, process topology, request lifecycle, technology rationale.\\
3 & Authentication, Authorization, Security \& Data Persistence & Trusted-proxy identity, RBAC, hardening, and the full data model.\\
4 & Problems, Submissions \& Code Execution Engine & Problem lifecycle, the judging pipeline, the sandbox, and code wrapping.\\
5 & Contest \& Assessment Engine & The contest lifecycle, deferred grading, auto-submission, feedback gate, and standings.\\
6 & Proctoring \& Academic Integrity & Client enforcement, server-side scoring, and integrity guarantees.\\
7 & Deployment, Infrastructure, Operations \& Future Scope & Topology, configuration, scaling, runbooks, and the roadmap.\\
\bottomrule
\end{longtable}

\vspace{0.6cm}
\begin{center}
\textit{--- End of Module 1. Continued in Module 2: System Architecture \& Technology Stack. ---}
\end{center}

\end{document}
